Impracticability requires more than a missing historical workbook. Under
AASB~108, the entity must have made every reasonable effort to apply the
requirement. Relevant difficulties include effects that cannot be determined,
assumptions that would have to be made about management's past intentions,
and significant estimates for which the necessary contemporaneous evidence
cannot be objectively distinguished from other information. Those conditions
need to be established for the reconstruction at issue
\citep[paragraph~5]{aasb108comp2022}.

An estimate prepared later can nevertheless describe an earlier period. The
question is whether it reflects the circumstances that existed then, using
evidence that would have been available when the earlier financial statements
were authorised for issue. A newly assembled estimate can meet that test;
replacing an earlier expectation with knowledge of what subsequently happened
uses hindsight. The date of preparation alone therefore does not settle the
question. Nor does authorisation establish when an outside investor could
obtain the information: that requires the separate public-availability record.
The insurance standard's specific transitional provisions govern their stated
cases, including permitted information under the modified and fair-value
routes described below
\citep[paragraphs~7, 19(a) and 50--53]{aasb108comp2022}.
